% GENERATED FILE. Edit canonical lessons, then run npm run generate:books.
% canonical-source-hash: fnv1a64:ce257d6af658e276
% canonical-lessons: MR-C288-ghatna, MR-C288-apghat, MR-C288-cori, MR-C288-budne, MR-C288-kirkol

\chapter{Something Happened}
\label{ch:mr-something-happened}
\hlchaptermodality{288}

\begin{chapteropening}
\textbf{By the end of this chapter:} \emph{I can say that something happened: an incident, an accident, a theft.}

Complete the last lesson of chapter 288: I can say that something happened: an incident, an accident, a theft.
\end{chapteropening}

\section[\texorpdfstring{\mr{घटना} (ghaṭnā)}{घटना (ghaṭnā)}]{\mr{घटना} (ghaṭnā) — an event; an incident}

\label{lesson:MR-C288-ghatna}



\begin{quote}
Before the new one: say the Marathi for dance, then the Marathi for a stage.
\end{quote}

\subsection*{You'll want to know: \mr{घटना}}
\textbf{\mr{घटना}} — \emph{ghaṭnā} — \textquotedblleft{}an event; an incident\textquotedblright{}.

\begin{quote}
\mr{ही} \mr{घटना} \mr{किती} \mr{मोठी} \mr{होती}, \mr{हे} \mr{तुम्ही} \mr{कोणाला} \mr{विचारता} \mr{यावर} \mr{अवलंबून} \mr{आहे}.
\end{quote}

\begin{quote}
\emph{hī ghaṭnā kitī moṭhī hotī, he tumhī koṇālā vicārtā yāvar avlambūn āhe.}
\end{quote}

\begin{quote}
How big this incident was depends on whom you ask.
\end{quote}

Feminine.

\mr{घटना} \mr{घडणे}, an event takes place; the oblique is \mr{घटने}-: \mr{घटनेचा}, of the event.

\begin{cousinweb}
\textbf{\mr{घटना}} (\emph{ghaṭnā}) is a tatsama, taken from Sanskrit \emph{ghaṭanā}, 'exertion; a taking effect, a coming about' (Monier-Williams). Marathi uses it for anything that happens.
\end{cousinweb}

\begin{grammarlens}[title={What you've built}]
One more word for telling what happened.
\end{grammarlens}

\subsection*{Your turn}
\begin{itemize}
\raggedright
  \item \emph{Say it:} \emph{ghaṭnā}
  \item \emph{Say it:} \emph{ghaṭnā}, once more
  \item \emph{Say it:} \emph{hī ghaṭnā kitī moṭhī hotī, he tumhī koṇālā vicārtā yāvar avlambūn āhe.} — How big this incident was depends on whom you ask.
  \item \emph{Recall:} say the Marathi for dance, then the Marathi for a stage, then say \emph{ghaṭnā} again
\end{itemize}

\begin{tcolorbox}[breakable,colback=teal!4,colframe=teal!35!black,title={Before you move on}]
What does \mr{घटना} mean? (\textquotedblleft{}An event; an incident\textquotedblright{}.) Say it once more.
\end{tcolorbox}
\section[\texorpdfstring{\mr{अपघात} (apghāt)}{अपघात (apghāt)}]{\mr{अपघात} (apghāt) — an accident}

\label{lesson:MR-C288-apghat}



\begin{quote}
Before the new one: say the Marathi for a stage, then the Marathi for an event.
\end{quote}

\subsection*{You'll want to know: \mr{अपघात}}
\textbf{\mr{अपघात}} — \emph{apghāt} — \textquotedblleft{}an accident\textquotedblright{}.

\begin{quote}
\mr{शाळेजवळ} \mr{काल} \mr{एक} \mr{मोठा} \mr{अपघात} \mr{झाला}.
\end{quote}

\begin{quote}
\emph{śāḷejvaḷ kāl ek moṭhā apghāt jhālā.}
\end{quote}

\begin{quote}
There was a big accident near the school yesterday.
\end{quote}

Masculine.

\mr{अपघात} \mr{होणे}, an accident happens; \mr{अपघातात}, in an accident.

\begin{cousinweb}
\textbf{\mr{अपघात}} (\emph{apghāt}) is a tatsama, taken from Sanskrit \emph{apaghāta}, 'a striking off, a warding off' (Monier-Williams); the sense 'an accident' is the modern one.
\end{cousinweb}

\begin{grammarlens}[title={What you've built}]
One more word for telling what happened.
\end{grammarlens}

\subsection*{Your turn}
\begin{itemize}
\raggedright
  \item \emph{Say it:} \emph{apghāt}
  \item \emph{Say it:} \emph{apghāt}, once more
  \item \emph{Say it:} \emph{śāḷejvaḷ kāl ek moṭhā apghāt jhālā.} — There was a big accident near the school yesterday.
  \item \emph{Recall:} say the Marathi for a stage, then the Marathi for an event, then say \emph{apghāt} again
\end{itemize}

\begin{tcolorbox}[breakable,colback=teal!4,colframe=teal!35!black,title={Before you move on}]
What does \mr{अपघात} mean? (\textquotedblleft{}An accident\textquotedblright{}.) Say it once more.
\end{tcolorbox}
\section[\texorpdfstring{\mr{चोरी} (corī)}{चोरी (corī)}]{\mr{चोरी} (corī) — a theft}

\label{lesson:MR-C288-cori}



\begin{quote}
Before the new one: say the Marathi for an event, then the Marathi for an accident.
\end{quote}

\subsection*{You'll want to know: \mr{चोरी}}
\textbf{\mr{चोरी}} — \emph{corī} — \textquotedblleft{}a theft\textquotedblright{}.

\begin{quote}
\mr{आमच्या} \mr{गल्लीत} \mr{दोन} \mr{घरांमध्ये} \mr{चोरी} \mr{झाली}.
\end{quote}

\begin{quote}
\emph{āmacyā gallīt don gharāmmadhye corī jhālī.}
\end{quote}

\begin{quote}
There was a theft in two houses on our lane.
\end{quote}

Feminine.

\mr{चोरी} \mr{होणे}, a theft happens; the verb \mr{चोरणे} is to steal.

\begin{cousinweb}
\textbf{\mr{चोरी}} (\emph{corī}) is a tadbhava, inherited from Sanskrit \emph{caurikā}, 'theft' (CDIAL 4937), which is built on \emph{cora}, 'a thief'.
\end{cousinweb}

\begin{grammarlens}[title={What you've built}]
One more word for telling what happened.
\end{grammarlens}

\subsection*{Your turn}
\begin{itemize}
\raggedright
  \item \emph{Say it:} \emph{corī}
  \item \emph{Say it:} \emph{corī}, once more
  \item \emph{Say it:} \emph{āmacyā gallīt don gharāmmadhye corī jhālī.} — There was a theft in two houses on our lane.
  \item \emph{Recall:} say the Marathi for an event, then the Marathi for an accident, then say \emph{corī} again
\end{itemize}

\begin{tcolorbox}[breakable,colback=teal!4,colframe=teal!35!black,title={Before you move on}]
What does \mr{चोरी} mean? (\textquotedblleft{}A theft\textquotedblright{}.) Say it once more.
\end{tcolorbox}
\section[\texorpdfstring{\mr{बुडणे} (buḍṇe)}{बुडणे (buḍṇe)}]{\mr{बुडणे} (buḍṇe) — to sink; to drown}

\label{lesson:MR-C288-budne}



\begin{quote}
Before the new one: say the Marathi for an accident, then the Marathi for a theft.
\end{quote}

\subsection*{You'll want to know: \mr{बुडणे}}
\textbf{\mr{बुडणे}} — \emph{buḍṇe} — \textquotedblleft{}to sink; to drown\textquotedblright{}.

\begin{quote}
\mr{वादळात} \mr{एक} \mr{जुनी} \mr{होडी} \mr{नदीत} \mr{बुडाली}.
\end{quote}

\begin{quote}
\emph{vādḷāt ek junī hoḍī nadīt buḍālī.}
\end{quote}

\begin{quote}
An old boat sank in the river during the storm.
\end{quote}

Past (he …): \mr{तो} \mr{बुडाला}, \emph{to buḍālā}.

The water takes -\mr{त}: \mr{पाण्यात} \mr{बुडणे}, to sink in the water.

\begin{cousinweb}
\textbf{\mr{बुडणे}} (\emph{buḍṇe}) is inherited: CDIAL lists Marathi \emph{buḍṇe}, 'to be drowned', under \emph{*buḍyati}, 'sinks' (CDIAL 9272).
\end{cousinweb}

\begin{grammarlens}[title={What you've built}]
One more word for telling what happened.
\end{grammarlens}

\subsection*{Your turn}
\begin{itemize}
\raggedright
  \item \emph{Say it:} \emph{buḍṇe}
  \item \emph{Say it:} \emph{buḍṇe}, once more
  \item \emph{Say it:} \emph{vādḷāt ek junī hoḍī nadīt buḍālī.} — An old boat sank in the river during the storm.
  \item \emph{Recall:} say the Marathi for an accident, then the Marathi for a theft, then say \emph{buḍṇe} again
\end{itemize}

\begin{tcolorbox}[breakable,colback=teal!4,colframe=teal!35!black,title={Before you move on}]
What does \mr{बुडणे} mean? (\textquotedblleft{}To sink; to drown\textquotedblright{}.) Say it once more.
\end{tcolorbox}
\section[\texorpdfstring{\mr{किरकोळ} (kirkoḷ)}{किरकोळ (kirkoḷ)}]{\mr{किरकोळ} (kirkoḷ) — minor; petty}

\label{lesson:MR-C288-kirkol}



\begin{quote}
Before the new one: say the Marathi for a theft, then the Marathi for to sink.
\end{quote}

\subsection*{You'll want to know: \mr{किरकोळ}}
\textbf{\mr{किरकोळ}} — \emph{kirkoḷ} — \textquotedblleft{}minor; petty\textquotedblright{}.

\begin{quote}
\mr{अपघात} \mr{किरकोळ} \mr{होता}, \mr{कोणालाही} \mr{लागलं} \mr{नाही}.
\end{quote}

\begin{quote}
\emph{apghāt kirkoḷ hotā, koṇālāhī lāglaṁ nāhī.}
\end{quote}

\begin{quote}
The accident was minor; nobody was hurt.
\end{quote}

\mr{किरकोळ} keeps one form: \mr{किरकोळ} \mr{अपघात}, \mr{किरकोळ} \mr{जखम}. It also means retail: \mr{किरकोळ} \mr{दुकान}, a retail shop.

\begin{cousinweb}
\textbf{\mr{किरकोळ}} (\emph{kirkoḷ}) is a Marathi word for small, scattered things, from retail goods to minor injuries; its origin is not settled.
\end{cousinweb}

\begin{grammarlens}[title={What you've built}]
One more word for telling what happened.
\end{grammarlens}

\subsection*{Your turn}
\begin{itemize}
\raggedright
  \item \emph{Say it:} \emph{kirkoḷ}
  \item \emph{Say it:} \emph{kirkoḷ}, once more
  \item \emph{Say it:} \emph{apghāt kirkoḷ hotā, koṇālāhī lāglaṁ nāhī.} — The accident was minor; nobody was hurt.
  \item \emph{Recall:} say the Marathi for a theft, then the Marathi for to sink, then say \emph{kirkoḷ} again
\end{itemize}

\begin{tcolorbox}[breakable,colback=teal!4,colframe=teal!35!black,title={Before you move on}]
What does \mr{किरकोळ} mean? (\textquotedblleft{}Minor; petty\textquotedblright{}.) Say it once more.
\end{tcolorbox}
